\documentclass[11pt]{article}

\usepackage[margin=1in]{geometry}
\usepackage{amsmath, amssymb}
\usepackage{listings}
\usepackage{xcolor}
\usepackage{booktabs}
\usepackage{hyperref}
\usepackage{array}
\usepackage{mdframed}

\definecolor{codebg}{RGB}{245,245,245}
\definecolor{codekw}{RGB}{0,90,180}
\definecolor{codecm}{RGB}{110,110,110}
\definecolor{codestr}{RGB}{160,40,40}
\definecolor{notebg}{RGB}{255,248,230}
\definecolor{noteborder}{RGB}{200,150,40}

\lstset{
  language=Python,
  backgroundcolor=\color{codebg},
  basicstyle=\ttfamily\small,
  keywordstyle=\color{codekw}\bfseries,
  commentstyle=\color{codecm}\itshape,
  stringstyle=\color{codestr},
  showstringspaces=false,
  breaklines=true,
  frame=single,
  numbers=left,
  numberstyle=\tiny\color{gray},
  numbersep=8pt,
  xleftmargin=18pt,
  tabsize=4,
}

\newmdenv[backgroundcolor=notebg, linecolor=noteborder, linewidth=1pt,
          innertopmargin=6pt, innerbottommargin=6pt, innerleftmargin=8pt, innerrightmargin=8pt]{notebox}

\newcommand{\eq}[1]{\begin{equation}#1\end{equation}}

\title{Line-by-Line Companion to \texttt{fiscal\_model.py}:\\
Every Line of Code and Its Mathematics}
\author{}
\date{}

\begin{document}
\maketitle

\begin{abstract}
\noindent This document explains \emph{every} line of \texttt{fiscal\_model.py}
--- the model behind the Streamlit app that replicates and extends Baxter \&
King (1993, \textit{AER}), ``Fiscal Policy in General Equilibrium'' ---
pairing each statement with the exact mathematical object it represents,
cited as \textbf{L}$n$ for line $n$. The model goes beyond the paper's
baseline in one respect the app exposes as a first-class control:
\textbf{productive public capital} (the paper's Section VI extension), with
production $Y_t=AK_t^{\theta_K}(K^G_t)^{\theta_G}N_t^{\theta_N}$. There is no
separate ``basic purchases'' category (it was removed, since the utility
function never values government consumption directly): total government
spending splits into just public investment $I^G$ (resource-using,
accumulates into $K^G$) and lump-sum transfers $G_T$ (resource-neutral).
Where the code implements a nontrivial derivation (the log-linearization of
the household's Euler equation and labor-leisure condition, and the
Blanchard--Kahn/King--Plosser--Rebelo saddle-path solution via
eigendecomposition), that derivation is carried out here from the model's
first-order conditions, term by term. Section~\ref{sec:table4} additionally
documents \texttt{public\_investment\_long\_run}'s replication of the
paper's actual Table 4 (verified line-for-line against the published table),
including a genuine numerical instability that this function must avoid
(public capital is an \emph{essential} input whenever $\theta_G>0$, so
naively perturbing $I^G$ near $0$ makes the finite difference blow up).
\end{abstract}

\tableofcontents

\section{Module docstring and the model in one page (L1--L55)}

The docstring (L1--L46) states the model compactly:
\begin{align}
u_t &= \log C_t + \theta_L\log L_t, \qquad L_t = 1-N_t \tag{preferences}\\
Y_t &= A\,K_t^{\theta_K}(K^G_t)^{\theta_G}N_t^{\theta_N}, \qquad \theta_K=1-\theta_N \tag{technology}\\
K_{t+1} &= (1-\delta)K_t+I_t, \qquad K^G_{t+1}=(1-\delta)K^G_t+I^G_t \tag{capital accum.}\\
Y_t &= C_t+I_t+I^G_t \tag{resource constraint}\\
\tau_t Y_t &= I^G_t+G_{T,t}, \qquad \tau_t=\frac{I^G_t+G_{T,t}}{Y_t} \tag{government budget}
\end{align}
with two financing rules for a change in total government spending: both set
$\tau_t$ by the \emph{same} formula above; they differ only in whether that
tax is a genuine wedge on the household's factor income
(\textbf{income\_tax}, the paper's original ``distortionary''/GRH channel)
or leaves it untouched (\textbf{lump\_sum}, a non-distorting poll tax for
the same revenue). Two solution objects are provided: an exact closed-form
\texttt{steady\_state} (Section~\ref{sec:ss}) and a log-linearized
\texttt{simulate\_transition} (Sections \ref{sec:loglinear}--\ref{sec:transition}).
All quantities are \emph{great ratios} --- shares of output/the time
endowment --- so the model is scale-free (L43--45).

\begin{lstlisting}[firstnumber=48]
from __future__ import annotations
from dataclasses import dataclass
from typing import Literal, Optional
import numpy as np
Financing = Literal["lump_sum", "income_tax"]
\end{lstlisting}
\begin{itemize}
\item[\textbf{L48}] Postpones type-annotation evaluation (a typing
convenience only).
\item[\textbf{L50}] Imports \texttt{dataclass}, used for \texttt{SteadyState},
\texttt{TransitionPath}, \texttt{PolicyExperiment}.
\item[\textbf{L51}] Imports \texttt{Literal} (used at L55 to restrict the
\texttt{financing} argument to exactly two string values) and
\texttt{Optional} (used for \texttt{PolicyExperiment} fields that can be
\texttt{None} before a path is computed).
\item[\textbf{L53}] NumPy --- the only numerical dependency; the saddle-path
solver (Section~\ref{sec:saddle}) uses \texttt{numpy.linalg.eig} for a
$2\times2$ eigendecomposition, the only linear-algebra call in the whole
project.
\item[\textbf{L55}] Defines \texttt{Financing} as a type restricted to the
two literal strings \texttt{"lump\_sum"} and \texttt{"income\_tax"}.
\end{itemize}

\section{\texttt{\_supply\_side\_impl}: the price system with public capital (L62--L98)}
\label{sec:supply}

\begin{lstlisting}[firstnumber=62]
def _supply_side_impl(theta_N: float, delta: float, r: float, A: float, tau: float,
                       theta_G: float, s_IG: float, distortionary: bool) -> dict:
\end{lstlisting}
\begin{itemize}
\item[\textbf{L62--L63}] Signature: seven primitives (five original plus
$\theta_G$, the public-capital productivity parameter, and $s_{I^G}$, the
public-investment share of output) plus a financing flag, out comes a
dictionary of derived prices/ratios. Because production is constant-returns
in $(K,N)$ for any \emph{given} level of public capital $K^G$, the
capital/labor ratio $\kappa=K/N$ is pinned down by the capital-Euler
condition alone, without yet knowing labor supply $N$ (solved separately in
Section~\ref{sec:calib}).
\end{itemize}

\begin{lstlisting}[firstnumber=64]
theta_K = 1.0 - theta_N
tax_factor = (1.0 - tau) if distortionary else 1.0
if tax_factor <= 0:
    raise ValueError("Invalid parameters: (1-tau) must be positive under Income Tax financing.")

kappa = ((tax_factor * theta_K * A) / (r + delta)) ** (1.0 / (1.0 - theta_K))
alpha = A * kappa ** theta_K  # will be rescaled by KG^theta_G below
KG = 0.0
\end{lstlisting}
\begin{itemize}
\item[\textbf{L64}] $\theta_K=1-\theta_N$, capital's Cobb--Douglas share.
\item[\textbf{L65}] The tax wedge on the household's factor income:
$(1-\tau)$ under \texttt{income\_tax} financing, $1$ under \texttt{lump\_sum}
financing. This single scalar is \emph{the} difference between the two
financing rules everywhere in the file --- both set the same $\tau$
(Section~\ref{sec:ss}), but only \texttt{income\_tax} lets it enter as a
factor-income wedge here.
\item[\textbf{L66--L67}] Guards against a tax rate $\ge100\%$ under
income-tax financing.
\item[\textbf{L69}] Solves the household's after-tax capital-Euler condition
for the capital/labor ratio, \emph{ignoring public capital's contribution
for now} (corrected in the fixed-point loop below when $\theta_G>0$):
\eq{ (1-\tau)\,\theta_K\, A\,\kappa^{\theta_K-1} = r+\delta \label{eq:euler-ss0} }
(or without the $(1-\tau)$ factor under lump-sum). Solving for $\kappa$:
$\kappa = \big(\text{tax\_factor}\cdot\theta_K A/(r+\delta)\big)^{1/(1-\theta_K)}$.
\item[\textbf{L70}] Average product of labor \emph{holding $K^G$ at its
$\theta_G=0$ value}, $\alpha\equiv Y/N = A\kappa^{\theta_K}$; gets rescaled
below when public capital is present and productive.
\item[\textbf{L71}] Initializes $K^G=0$, the value used unless one of the
two branches below overrides it.
\end{itemize}

\begin{lstlisting}[firstnumber=72]
if theta_G > 0 and s_IG > 0:
    alpha_ppN = A * kappa ** theta_K  # Y/N ignoring KG, will multiply by KG^theta_G
    alpha_guess = alpha_ppN
    for _ in range(200):
        KG_ppN = s_IG * alpha_guess / delta  # KG per unit of N
        base = (tax_factor * theta_K * A * KG_ppN ** theta_G) / (r + delta)
        kappa = base ** (1.0 / (1.0 - theta_K))
        alpha_new = A * kappa ** theta_K * KG_ppN ** theta_G
        if abs(alpha_new - alpha_guess) < 1e-12:
            alpha_guess = alpha_new
            break
        alpha_guess = alpha_new
    alpha = alpha_guess
    KG = s_IG * alpha / delta  # KG per unit of N
elif s_IG > 0:
    KG = s_IG * alpha / delta
\end{lstlisting}
\label{sec:kgfixedpoint}
\begin{itemize}
\item[\textbf{L72}] Only run the (more expensive) fixed-point iteration when
public capital is actually productive ($\theta_G>0$) \emph{and} present
($s_{I^G}>0$); when $\theta_G=0$ (L90--91, see the box below) the physical
stock is still computed, just without any feedback into $\kappa,\alpha$.
\item[\textbf{Why a fixed point is needed when $\theta_G>0$.}] The
capital-Euler condition generalizes to
\eq{ \text{tax\_factor}\cdot\theta_K\,A\,(K^G/N)^{\theta_G}\,\kappa^{\theta_K-1} = r+\delta, \label{eq:euler-ss-kg}}
but $K^G$ is itself pinned down by $K^G=I^G/\delta$ (steady-state
accumulation identity), and $I^G=s_{I^G}Y$, and $Y$ depends on $\kappa$
through \eqref{eq:euler-ss-kg} --- circular, solved by iterating on
$\alpha\equiv Y/N$ until it stops moving. Quantities are kept \emph{per unit
of $N$} (\texttt{\_ppN} suffix) since $N$ is not yet known at this stage.
\item[\textbf{L73--L82}] The fixed-point sweep: L74--75 initialize the
guess; L76 up to $200$ iterations; L77 recomputes $K^G/N$; L78 recomputes
$\kappa$ from \eqref{eq:euler-ss-kg}; L79 recomputes
$\alpha_{\text{new}}=A\kappa^{\theta_K}(K^G/N)^{\theta_G}$; L80--82
convergence check.
\item[\textbf{L83--L84}] Finalizes $\alpha$ and $K^G$ (per unit of $N$) at
their fixed-point values.
\end{itemize}

\begin{notebox}
\textbf{L85--L86: public capital must still accumulate at $\theta_G=0$.} An
earlier version of this function set $K^G\equiv0$ whenever $\theta_G\le0$,
conflating ``not productive'' with ``does not exist.'' But $K^G$'s law of
motion, $K^G_{t+1}=(1-\delta)K^G_t+I^G_t$, does not care whether $\theta_G$
is $0$ --- at a steady state $K^G=I^G/\delta$ regardless, so a positive
public-investment share $s_{I^G}$ should always produce a nonzero physical
capital stock, even when that stock has no effect on production. This bug
was visible in the app: with $\theta_G=0$ and a large public-investment
composition share, the ``Public capital $K^G$'' row of the steady-state
table stayed at $0.000000$ while ``Public investment $I^G$'' was clearly
positive. The \texttt{elif s\_IG > 0} branch (L85--86) fixes this: no fixed
point is needed here (unlike the $\theta_G>0$ branch above) because
$(K^G)^0=1$ means $K^G$ cannot feed back into $\alpha,\kappa$ at all when
$\theta_G=0$, so $KG=s_{I^G}\alpha/\delta$ can be computed directly from the
already-known $\alpha$.
\end{notebox}

\begin{lstlisting}[firstnumber=95]
w = theta_N * alpha
q = tax_factor * theta_K * alpha / kappa
s_I = delta * kappa / alpha
return dict(theta_K=theta_K, kappa=kappa, alpha=alpha, w=w, q=q, s_I=s_I, KG=KG, tax_factor=tax_factor)
\end{lstlisting}
\begin{itemize}
\item[\textbf{L95}] The pre-tax real wage, $w=\theta_N\alpha=\theta_N\,Y/N$.
\item[\textbf{L96}] The after-tax rental rate,
$q=\text{tax\_factor}\cdot\theta_K\alpha/\kappa=\text{tax\_factor}\cdot\text{MPK}$
--- by construction (since $\kappa$ was solved from
\eqref{eq:euler-ss-kg}), $q\equiv r+\delta$ exactly.
\item[\textbf{L97}] $s_I\equiv I/Y=\delta K/Y=\delta\kappa/\alpha$.
\item[\textbf{L98}] Packages eight derived quantities, including
\texttt{tax\_factor} itself (reused by every caller) and \texttt{KG}
(public capital per unit of $N$).
\end{itemize}

\section{\texttt{SteadyState} and its constructors (L105--L200)}
\label{sec:ss}

\begin{lstlisting}[firstnumber=105]
@dataclass
class SteadyState:
    theta_N: float
    theta_K: float
    theta_G: float
    delta: float
    r: float
    A: float
    tau: float
    theta_L: float
    s_G: float       # total government spending share (IG+G_T)/Y
    s_IG: float
    s_GT: float
    distortionary: bool
    kappa: float
    alpha: float
    w: float
    q: float
    s_I: float
    s_C: float
    N: float
    L: float
    Y: float
    K: float
    KG: float
    C: float
    I: float
    G: float          # total government spending IG+G_T
    IG: float
    GT: float

    def as_dict(self) -> dict:
        return self.__dict__.copy()
\end{lstlisting}
\begin{itemize}
\item[\textbf{L105--L136}] A plain dataclass holding every parameter
\emph{and} every derived great ratio and level at one steady state: six
primitives, the calibrated $\theta_L$, the two-way composition
($s_G,s_{IG},s_{GT}$) and the financing flag, the seven
\texttt{\_supply\_side\_impl} outputs ($\kappa,\alpha,w,q,s_I,s_C$ --- $s_C$
is added here, not returned by that function itself), and the levels
($N,L,Y,K,K^G,C,I,G,I^G,G_T$).
\item[\textbf{L136--L137}] \texttt{as\_dict} returns a shallow copy of the
attribute dictionary --- used by \texttt{test\_model.py} to print every
field by name.
\end{itemize}

\begin{lstlisting}[firstnumber=140]
def calibrate_theta_L(theta_N: float, delta: float, r: float, A: float, theta_G: float,
                       s_IG: float, s_GT: float, distortionary: bool,
                       N_target: float) -> float:
    tau = s_IG + s_GT
    supply = _supply_side_impl(theta_N, delta, r, A, tau, theta_G, s_IG, distortionary)
    s_C = 1.0 - supply["s_I"] - s_IG
    if s_C <= 0:
        raise ValueError("Government share too large: steady-state consumption share <= 0.")
    N = N_target
    theta_L = supply["tax_factor"] * theta_N * (1.0 - N) / (s_C * N)
    if theta_L <= 0:
        raise ValueError("Implied theta_L <= 0; adjust N_target or the tax/spending settings.")
    return theta_L
\end{lstlisting}
\label{sec:calib}
\begin{itemize}
\item[\textbf{L140--L142}] Signature: given the technology/policy primitives
and a target labor supply $\bar N\equiv$\texttt{N\_target}, returns the
leisure-preference weight $\theta_L$ that makes steady-state labor supply
equal $\bar N$.
\item[\textbf{L144}] $\tau=s_{IG}+s_{GT}$ --- the government-budget
identity, holding regardless of financing rule.
\item[\textbf{L146}] Calls Section~\ref{sec:supply}'s function.
\item[\textbf{L147}] $s_C=1-s_I-s_{IG}$ --- note only $s_{IG}$ (not
$s_{GT}$) reduces the resource constraint, since transfers are rebated to
households rather than spent on goods.
\item[\textbf{L148--L149}] Guards against $s_C\le0$.
\item[\textbf{L151}] The calibration formula, from the labor-leisure FOC
\eq{ \theta_L\,\frac{C}{L} = \text{tax\_factor}\cdot w \label{eq:foc-labor} }
solved for $\theta_L$ using $w=\theta_N\alpha$, $C=s_C Y$, $\alpha=Y/N$,
$L=1-N$:
\eq{ \theta_L = \frac{\text{tax\_factor}\cdot\theta_N(1-N)}{s_C N}. }
\item[\textbf{L152--L153}] Guards against a non-positive $\theta_L$.
\end{itemize}

\begin{lstlisting}[firstnumber=157]
def steady_state(theta_N: float, delta: float, r: float, A: float, theta_G: float,
                  s_IG: float, s_GT: float, distortionary: bool,
                  theta_L: float) -> SteadyState:
    tau = s_IG + s_GT
    s_G = tau
    supply = _supply_side_impl(theta_N, delta, r, A, tau, theta_G, s_IG, distortionary)
    s_I = supply["s_I"]
    s_C = 1.0 - s_I - s_IG
    if s_C <= 0:
        raise ValueError("Government share too large: steady-state consumption share <= 0.")
\end{lstlisting}
\begin{itemize}
\item[\textbf{L157--L159}] Signature: same primitives plus $\theta_L$ (now
given, not calibrated), returning a full \texttt{SteadyState}.
\item[\textbf{L162--L168}] Re-derives $\tau,s_G,s_I,s_C$ as in
Section~\ref{sec:calib}.
\end{itemize}

\begin{lstlisting}[firstnumber=170]
tax_factor = supply["tax_factor"]
N = tax_factor * theta_N / (theta_L * s_C + tax_factor * theta_N)
if not (0 < N < 1):
    raise ValueError(f"Implied labor supply N={N:.3f} is outside (0,1).")
L = 1.0 - N
\end{lstlisting}
\begin{itemize}
\item[\textbf{L171}] Solves \eqref{eq:foc-labor} for $N$ directly: from
$\theta_L s_C N=\text{tax\_factor}\cdot\theta_N(1-N)$,
\eq{ N=\frac{\text{tax\_factor}\cdot\theta_N}{\theta_L s_C+\text{tax\_factor}\cdot\theta_N}. }
The algebraic inverse of L151 --- plugging L151's $\theta_L$ back in
returns exactly $N=$\texttt{N\_target}.
\item[\textbf{L172--L173}] Guards against infeasible labor supply.
\item[\textbf{L174}] $L=1-N$.
\end{itemize}

\begin{lstlisting}[firstnumber=176]
Y = supply["alpha"] * N
K = supply["kappa"] * N
KG = supply["KG"] * N
C = s_C * Y
I = s_I * Y
IG = s_IG * Y
GT = s_GT * Y
G = IG + GT
\end{lstlisting}
\begin{itemize}
\item[\textbf{L176--L178}] $Y=\alpha N$, $K=\kappa N$, and $K^G=(K^G/N)\cdot
N$ --- converting the per-$N$ ratios from Section~\ref{sec:kgfixedpoint}
into true levels now that $N$ is known.
\item[\textbf{L179--L180}] $C=s_C Y$, $I=s_I Y$.
\item[\textbf{L181--L182}] $I^G=s_{IG}Y$, $G_T=s_{GT}Y$.
\item[\textbf{L183}] $G=I^G+G_T$, total government spending; by
construction $C+I+I^G=(s_C+s_I+s_{IG})Y=Y$ exactly.
\end{itemize}

\begin{lstlisting}[firstnumber=185]
return SteadyState(
    theta_N=theta_N, theta_K=supply["theta_K"], theta_G=theta_G, delta=delta, r=r, A=A,
    tau=tau, theta_L=theta_L, s_G=s_G, s_IG=s_IG, s_GT=s_GT,
    distortionary=distortionary, kappa=supply["kappa"], alpha=supply["alpha"],
    w=supply["w"], q=supply["q"], s_I=s_I, s_C=s_C, N=N, L=L, Y=Y, K=K, KG=KG, C=C, I=I,
    G=G, IG=IG, GT=GT,
)
\end{lstlisting}
\begin{itemize}
\item[\textbf{L185--L191}] Assembles all twenty-eight fields into the
returned \texttt{SteadyState}.
\end{itemize}

\begin{lstlisting}[firstnumber=194]
def steady_state_for_policy(theta_N: float, delta: float, r: float, A: float, theta_G: float,
                             s_G_total: float, f_IG: float, f_GT: float,
                             theta_L: float, distortionary: bool) -> SteadyState:
    s_IG, s_GT = s_G_total * f_IG, s_G_total * f_GT
    return steady_state(theta_N, delta, r, A, theta_G, s_IG, s_GT, distortionary, theta_L)
\end{lstlisting}
\label{sec:policybudget}
\begin{itemize}
\item[\textbf{L194--L196}] Signature: takes a \emph{total}
government-spending ratio $s_{G,\text{total}}$ and fixed composition
fractions $f_{IG},f_{GT}$ (summing to $1$), rather than the two composition
levels directly. This is the function the app calls once for the baseline
policy and once for the counterfactual, holding the composition fractions
fixed across the two.
\item[\textbf{L199}] Splits the total by the given fractions.
\item[\textbf{L200}] Calls Section~\ref{sec:ss}'s \texttt{steady\_state}.
\item[\textbf{On the household budget and why $G_T$ drops out of the
resource constraint.}] The government budget is
$\tau_tY_t=I^G_t+G_{T,t}$. The household budget differs by financing rule:
\eq{
\text{income\_tax:}\quad C_t+I_t=(1-\tau_t)(w_tN_t+r_tK_t)+G_{T,t}
\qquad
\text{lump\_sum:}\quad C_t+I_t=w_tN_t+r_tK_t-\tau_tY_t+G_{T,t}.
}
Combined with $w_tN_t+r_tK_t=Y_t$ and the government budget identity, both
cases reduce \emph{in the aggregate} to $Y_t=C_t+I_t+I^G_t$ --- transfers
are received by the household but exactly offset by the tax it pays, so
they wash out of the resource constraint even though they matter for
household-level after-tax income.
\end{itemize}

\section{The log-linear coefficients: derivation and code (L207--L259)}
\label{sec:loglinear}

The code's own docstring (L208--L213) states the four blocks to be solved:
\begin{align}
\hat y_t &= y_k\hat k_t+y_{kg}\hat k^G_t+y_c\hat c_t+y_g\hat g_t &&\text{(static block)}\label{eq:static}\\
\hat k_{t+1} &= \Phi_{kk}\hat k_t+\Phi_{kc}\hat c_t+\Phi_{kg}\hat g_t+\Phi_{k,kg}\hat k^G_t &&\text{(capital accumulation)}\label{eq:capaccum}\\
\hat c_t &= \text{coef}_c\hat c_{t+1}+\text{coef}_k\hat k_{t+1}+\text{coef}_g\hat g_{t+1}+\text{coef}_{kg}\hat k^G_{t+1} &&\text{(Euler equation)}\label{eq:euler}
\end{align}
where a hat denotes a log-deviation from the steady state \texttt{ss}. This
section derives all twelve coefficients from the production function, the
labor-leisure FOC \eqref{eq:foc-labor}, and the household's capital-Euler
equation.

\begin{lstlisting}[firstnumber=207]
def _log_linear_coeffs(ss: SteadyState) -> dict:
\end{lstlisting}
\begin{itemize}
\item[\textbf{L207}] Signature: takes only the steady state to linearize
around; the financing flag lives on \texttt{ss.distortionary}.
\end{itemize}

\begin{lstlisting}[firstnumber=215]
theta_N, theta_K, theta_G = ss.theta_N, ss.theta_K, ss.theta_G
L = ss.L
tau = ss.tau
s_C, s_I, s_IG = ss.s_C, ss.s_I, ss.s_IG
delta, r = ss.delta, ss.r
\end{lstlisting}
\begin{itemize}
\item[\textbf{L215--L219}] Unpacks every steady-state quantity needed
below, including $s_{IG}$ (the resource-using share; $s_{GT}$ is not
needed here, since transfers never enter the resource-constraint-based
formulas that follow).
\end{itemize}

\subsection{The static block: output as a function of \texorpdfstring{$\hat k,\hat k^G,\hat c,\hat g$}{k,kG,c,g} (L221--L236)}

\textbf{Derivation.} Log-linearize the production function:
\eq{ \hat y_t = \theta_K\hat k_t+\theta_G\hat k^G_t+\theta_N\hat n_t. \label{eq:prod-loglin} }
Log-linearize the labor-leisure FOC \eqref{eq:foc-labor}: writing it as
$\theta_L C_t=\text{tax\_factor}_t\theta_N Y_t(1-N_t)/N_t$ and
differentiating, using $d\log(1-N_t)=-\frac{N}{L}\hat n_t$,
\eq{ \hat c_t = \widehat{(\text{tf})}_t+\hat y_t-\frac{\hat n_t}{L}. \label{eq:foc-loglin}}

\textbf{Case \texttt{lump\_sum} (tax\_factor$\equiv1$, $\widehat{(\text{tf})}_t\equiv0$).}
Solving \eqref{eq:foc-loglin} for $\hat n_t$: $\hat n_t=L(\hat y_t-\hat c_t)$.
Substituting into \eqref{eq:prod-loglin}:
\eq{ \hat y_t\underbrace{(1-\theta_N L)}_{\displaystyle D}
       = \theta_K\hat k_t+\theta_G\hat k^G_t-\theta_N L\,\hat c_t, }
\eq{ \Longrightarrow\quad y_k=\frac{\theta_K}{D}, \quad y_{kg}=\frac{\theta_G}{D}, \quad y_c=-\frac{\theta_N L}{D}, \quad y_g=0. \label{eq:static-lump}}

\begin{lstlisting}[firstnumber=221]
D = 1.0 - theta_N * L
if ss.distortionary:
    Dp = D - theta_N * L * tau / (1.0 - tau)
    y_k = theta_K / Dp
    y_kg = theta_G / Dp
    y_c = -theta_N * L / Dp
    y_g = -theta_N * L * tau / ((1.0 - tau) * Dp)
    tax_wedge = tau / (1.0 - tau)
else:
    y_k = theta_K / D
    y_kg = theta_G / D
    y_c = -theta_N * L / D
    y_g = 0.0
    tax_wedge = 0.0
\end{lstlisting}
\begin{itemize}
\item[\textbf{L221}] $D\equiv1-\theta_N L$.
\item[\textbf{L231--L235}] Exactly implements \eqref{eq:static-lump}: no
$\hat g_t$ term because $g$ never enters the lump-sum labor-leisure FOC,
and $y_{kg}=\theta_G/D$ shares $y_k$'s denominator since $\hat k^G_t$
entered \eqref{eq:prod-loglin} exactly like $\hat k_t$.
\end{itemize}

\textbf{Case \texttt{income\_tax}.} Differentiating $\tau_tY_t=G_t$
(total government spending) gives
\eq{ \widehat{(\text{tf})}_t=\frac{\tau}{1-\tau}(\hat y_t-\hat g_t). \label{eq:tauhat} }
Substituting into \eqref{eq:foc-loglin} and solving for $\hat n_t$:
\eq{ \hat n_t = L\left[\frac{\hat y_t}{1-\tau}-\hat c_t-\frac{\tau}{1-\tau}\hat g_t\right] \label{eq:n-distortionary} }
--- exactly the formula \texttt{simulate\_transition} uses directly
(Section~\ref{sec:transition}, L376--L379). Substituting back into
\eqref{eq:prod-loglin} and collecting $\hat y_t$ terms:
\eq{
\hat y_t\left(\underbrace{D-\frac{\theta_N L\tau}{1-\tau}}_{\displaystyle D_p}\right)
   = \theta_K\hat k_t+\theta_G\hat k^G_t-\theta_N L\hat c_t-\frac{\theta_N L\tau}{1-\tau}\hat g_t.
   \label{eq:static-distortionary-correct}
}
\begin{itemize}
\item[\textbf{L222}] $D_p\equiv D-\theta_N L\tau/(1-\tau)$.
\item[\textbf{L223--L226}] $y_k=\theta_K/D_p$, $y_{kg}=\theta_G/D_p$,
$y_c=-\theta_N L/D_p$, $y_g=-\theta_N L\tau/[(1-\tau)D_p]$ --- matching
\eqref{eq:static-distortionary-correct}.
\item[\textbf{L227}] $\text{tax\_wedge}\equiv\tau/(1-\tau)$, reused in the
Euler-equation coefficients (Section~\ref{sec:euler}).
\end{itemize}

\subsection{The Euler equation coefficients (L238--L250)}
\label{sec:euler}

\textbf{Derivation.} The after-tax capital-Euler equation is
\eq{ \frac{1}{C_t} = \beta\,E_t\!\left[\frac{1+\text{tf}_{t+1}\,R_{t+1}-\delta}{C_{t+1}}\right], \qquad
     R_{t+1}\equiv\theta_K\frac{Y_{t+1}}{K_{t+1}}. }
At steady state $1/\beta=\text{tf}\cdot R+1-\delta=q+1-\delta$, and $q\equiv
r+\delta$, so $\beta=1/(1+r)$. Writing $X_{t+1}\equiv\text{tf}_{t+1}R_{t+1}+1-\delta$
and log-linearizing $1/C_t=\beta X_{t+1}/C_{t+1}$:
\eq{ \hat c_t=\hat c_{t+1}-\hat X_{t+1}, \qquad
     \hat X_{t+1} = \underbrace{\beta q}_{\displaystyle\Omega}
     \Big[\widehat{(\text{tf})}_{t+1}+\hat R_{t+1}\Big]. }
Since $R_t=\theta_K Y_t/K_t$, $\hat R_t=\hat y_t-\hat k_t$; using the static
block, $\hat y_t-\hat k_t=(y_k-1)\hat k_t+y_{kg}\hat k^G_t+y_c\hat c_t+y_g\hat
g_t$. Substituting this and \eqref{eq:tauhat} (or $0$, lump-sum), then $\hat
y_{t+1}$'s own static-block expression where needed, and collecting
coefficients on $\hat k_{t+1},\hat k^G_{t+1},\hat c_{t+1},\hat g_{t+1}$ in
$\hat c_t=\hat c_{t+1}-\hat X_{t+1}$, gives the four coefficients below.

\begin{lstlisting}[firstnumber=238]
beta = 1.0 / (1.0 + r)
q = r + delta
Omega = beta * q  # = q/(1+r)

ymk_k = y_k - 1.0
ymk_c = y_c
ymk_g = y_g
ymk_kg = y_kg

coef_c = 1.0 - Omega * ymk_c - Omega * tax_wedge * y_c
coef_k = -Omega * ymk_k - Omega * tax_wedge * y_k
coef_g = -Omega * ymk_g + Omega * tax_wedge * (1.0 - y_g)
coef_kg = -Omega * ymk_kg - Omega * tax_wedge * y_kg
\end{lstlisting}
\begin{itemize}
\item[\textbf{L238}] $\beta=1/(1+r)$.
\item[\textbf{L239}] $q=r+\delta$.
\item[\textbf{L240}] $\Omega\equiv\beta q$.
\item[\textbf{L242--L245}] $\hat R_t=\hat y_t-\hat k_t$'s coefficients:
\texttt{ymk\_k}$=y_k-1$, \texttt{ymk\_c}$=y_c$, \texttt{ymk\_g}$=y_g$,
\texttt{ymk\_kg}$=y_{kg}$.
\item[\textbf{L247}] $\text{coef}_c=1-\Omega\cdot\text{ymk\_c}-\Omega\cdot\text{tax\_wedge}\cdot y_c$.
\item[\textbf{L248}] $\text{coef}_k=-\Omega\cdot\text{ymk\_k}-\Omega\cdot\text{tax\_wedge}\cdot y_k$.
\item[\textbf{L249}] $\text{coef}_g=-\Omega\cdot\text{ymk\_g}+\Omega\cdot\text{tax\_wedge}\cdot(1-y_g)$.
\item[\textbf{L250}] $\text{coef}_{kg}=-\Omega\cdot\text{ymk\_kg}-\Omega\cdot\text{tax\_wedge}\cdot y_{kg}$
--- structurally identical to \texttt{coef\_k}'s template (both $\hat
k_{t+1}$ and $\hat k^G_{t+1}$ affect the Euler equation only through $\hat
R_{t+1}=\hat y_{t+1}-\hat k_{t+1}$), minus the ``$-1$'' term since $\hat
k^G_{t+1}$ never appears \emph{subtracted} the way $\hat k_{t+1}$ does.
\end{itemize}

\subsection{The capital-accumulation coefficients (L252--L255)}
\label{sec:capaccumsec}

\textbf{Derivation.} Log-linearizing $K_{t+1}=(1-\delta)K_t+I_t$ (using
$I=\delta K$ at steady state) and the resource constraint $Y_t=C_t+I_t+I^G_t$
(using $s_C+s_I+s_{IG}=1$):
\eq{ \hat\imath_t = \frac{y_k}{s_I}\hat k_t+\frac{y_{kg}}{s_I}\hat k^G_t+\frac{y_c-s_C}{s_I}\hat c_t+\frac{y_g-s_{IG}}{s_I}\hat g_t, }
substituted into $\hat k_{t+1}=(1-\delta)\hat k_t+\delta\hat\imath_t$:

\begin{lstlisting}[firstnumber=252]
Phi_kk = (1.0 - delta) + delta * y_k / s_I
Phi_kc = delta * (y_c - s_C) / s_I
Phi_kg = delta * (y_g - s_IG) / s_I
Phi_k_kg = delta * y_kg / s_I
\end{lstlisting}
\begin{itemize}
\item[\textbf{L252}] $\Phi_{kk}=(1-\delta)+\delta y_k/s_I$.
\item[\textbf{L253}] $\Phi_{kc}=\delta(y_c-s_C)/s_I$.
\item[\textbf{L254}] $\Phi_{kg}=\delta(y_g-s_{IG})/s_I$ --- $s_{IG}$, not
$s_G$: only resource-using spending crowds out investment.
\item[\textbf{L255}] $\Phi_{k,kg}=\delta y_{kg}/s_I$ --- public capital
affects next period's private capital only through its effect on $\hat
y_t$, with no direct resource-crowding term (it's a stock, not a
competing flow).
\end{itemize}

\begin{lstlisting}[firstnumber=257]
return dict(y_k=y_k, y_kg=y_kg, y_c=y_c, y_g=y_g, s_GR=s_IG,
            coef_c=coef_c, coef_k=coef_k, coef_g=coef_g, coef_kg=coef_kg,
            Phi_kk=Phi_kk, Phi_kc=Phi_kc, Phi_kg=Phi_kg, Phi_k_kg=Phi_k_kg)
\end{lstlisting}
\begin{itemize}
\item[\textbf{L257--L258}] Returns all twelve coefficients plus $s_{IG}$
(keyed \texttt{"s\_GR"} for historical reasons, predating the removal of
the separate ``basic purchases'' category, when the resource-using share
was $s_{GB}+s_{IG}$ rather than just $s_{IG}$).
\end{itemize}

\section{\texttt{\_kg\_hat\_path}: public capital's own law of motion (L262--L273)}
\label{sec:kghat}

\textbf{Derivation.} Log-linearizing $K^G_{t+1}=(1-\delta)K^G_t+I^G_t$
(using $I^G=\delta K^G$ at steady state), and noting that public investment
moves proportionally with total government spending at fixed composition
fractions ($I^G_t=f_{IG}G_t\Rightarrow\hat\imath^G_t=\hat g_t$):
\eq{ \hat k^G_{t+1} = (1-\delta)\hat k^G_t+\delta\,\hat g_t. \label{eq:kg-law} }
Public capital's law of motion does not depend on $\hat k_t$ or $\hat c_t$
at all --- it is driven exogenously by the (known) government-spending path
--- which is why the model does not need a $3\times3$ generalization of the
Blanchard--Kahn eigendecomposition (Section~\ref{sec:saddle} stays
$2\times2$).

\begin{lstlisting}[firstnumber=262]
def _kg_hat_path(delta: float, g_full: np.ndarray, kg0: float) -> np.ndarray:
    T = len(g_full)
    kg = np.zeros(T + 1)
    kg[0] = kg0
    g_ext = np.concatenate([g_full, [0.0]])
    for t in range(T):
        kg[t + 1] = (1.0 - delta) * kg[t] + delta * g_ext[t]
    return kg
\end{lstlisting}
\begin{itemize}
\item[\textbf{L262}] Signature: returns $T+1$ values, $\hat k^G_0,\dots,\hat
k^G_T$ --- one more than the input length, so the Euler-equation forcing at
the last simulated period $t=T-1$ (needing $\hat k^G_T$) is always
available.
\item[\textbf{L267}] Seeds $\hat k^G_0=$\texttt{kg0}.
\item[\textbf{L270}] Appends a trailing zero (shock has died out by the
horizon's end).
\item[\textbf{L271--L272}] Implements \eqref{eq:kg-law} by simple forward
iteration --- safe here (a stable, first-order, one-variable recursion with
root $1-\delta\in(0,1)$).
\end{itemize}

\section{\texttt{\_solve\_saddle\_path}: the Blanchard--Kahn/KPR solution (L276--L337)}
\label{sec:saddle}

\begin{lstlisting}[firstnumber=276]
def _solve_saddle_path(Phi_kk: float, Phi_kc: float, Phi_kg: float, Phi_k_kg: float,
                        coef_c: float, coef_k: float, coef_g: float, coef_kg: float,
                        k0: float, g_full: np.ndarray, kg_full: np.ndarray) -> tuple[np.ndarray, np.ndarray]:
\end{lstlisting}
\begin{itemize}
\item[\textbf{L276--L278}] Signature: the eight coefficients from
Section~\ref{sec:loglinear} involving $\hat k$ or $\hat c$, an initial
capital log-deviation $k_0$, \texttt{g\_full}, and the length-$(T+1)$
public-capital path \texttt{kg\_full} from Section~\ref{sec:kghat},
supplied as a \emph{known} forcing input.
\end{itemize}

\textbf{Setting up a first-order forward system.} Solving \eqref{eq:euler}
for $\hat c_{t+1}$ using \eqref{eq:capaccum}'s expression for $\hat
k_{t+1}$ and substituting gives:
\eq{
\begin{pmatrix}\hat k_{t+1}\\ \hat c_{t+1}\end{pmatrix}
= \underbrace{\begin{pmatrix}\Phi_{kk} & \Phi_{kc}\\[2pt]
              -\dfrac{\text{coef}_k\Phi_{kk}}{\text{coef}_c} & \dfrac{1-\text{coef}_k\Phi_{kc}}{\text{coef}_c}
              \end{pmatrix}}_{\displaystyle M}
  \begin{pmatrix}\hat k_t\\ \hat c_t\end{pmatrix}
  +\begin{pmatrix}\Phi_{kg}\hat g_t+\Phi_{k,kg}\hat k^G_t\\[2pt]
                   -\dfrac{\text{coef}_k(\Phi_{kg}\hat g_t+\Phi_{k,kg}\hat k^G_t)}{\text{coef}_c}-\dfrac{\text{coef}_g}{\text{coef}_c}\hat g_{t+1}-\dfrac{\text{coef}_{kg}}{\text{coef}_c}\hat k^G_{t+1}\end{pmatrix}.
  \label{eq:M-system}
}
$M$ itself is unaffected by public capital --- built only from
$\Phi_{kk},\Phi_{kc},\text{coef}_k,\text{coef}_c$; public capital enters
purely through the additive forcing vector.

\begin{lstlisting}[firstnumber=298]
M = np.array([[Phi_kk, Phi_kc],
              [-coef_k * Phi_kk / coef_c, (1.0 - coef_k * Phi_kc) / coef_c]])
\end{lstlisting}
\begin{itemize}
\item[\textbf{L298--L299}] Builds $M$ exactly as in \eqref{eq:M-system}.
\end{itemize}

\textbf{Diagonalizing $M$.} $M=V\,\text{diag}(\lambda_s,\lambda_u)\,V^{-1}$;
saddle-path stability requires $|\lambda_s|<1<|\lambda_u|$.

\begin{lstlisting}[firstnumber=300]
vals, vecs = np.linalg.eig(M)
if np.max(np.abs(vals.imag)) > 1e-8:
    raise RuntimeError("Complex eigenvalues encountered; parameters imply oscillatory "
                        "dynamics not supported by this simplified solver.")
vals = vals.real
vecs = vecs.real
stable_idx = int(np.argmin(np.abs(vals)))
unstable_idx = 1 - stable_idx
lam_s, lam_u = vals[stable_idx], vals[unstable_idx]
if not (abs(lam_s) < 1.0 < abs(lam_u)):
    raise RuntimeError("Model is not saddle-path stable for these parameters "
                        "(need exactly one root inside and one outside the unit circle).")
V = vecs[:, [stable_idx, unstable_idx]]
Vinv = np.linalg.inv(V)
\end{lstlisting}
\begin{itemize}
\item[\textbf{L300}] NumPy's eigenvalue solver.
\item[\textbf{L301--L304}] Guards against complex eigenvalues; discards
negligible imaginary parts.
\item[\textbf{L306}] Stable eigenvalue = smaller absolute value.
\item[\textbf{L307}] The other index is unstable.
\item[\textbf{L308}] Reads off $\lambda_s,\lambda_u$.
\item[\textbf{L309--L311}] Verifies true saddle-path stability.
\item[\textbf{L312}] Reorders $V$'s columns (stable, unstable).
\item[\textbf{L313}] Inverts $V$, needed to rotate $(\hat k_t,\hat c_t)$
into the eigenbasis where the system decouples.
\end{itemize}

\begin{lstlisting}[firstnumber=315]
T = len(g_full)
g_ext = np.concatenate([g_full, [0.0]])
kg_now = kg_full[:T]     # kg_hat_t for t=0..T-1
kg_next = kg_full[1:T + 1]  # kg_hat_{t+1} for t=0..T-1
Nk = Phi_kg * g_ext[:T] + Phi_k_kg * kg_now
Nc = (-coef_k * (Phi_kg * g_ext[:T] + Phi_k_kg * kg_now) - coef_g * g_ext[1:T + 1]
      - coef_kg * kg_next) / coef_c
N = np.stack([Nk, Nc], axis=1)
n = N @ Vinv.T  # forcing rotated into the eigen-basis: n[:,0]=stable comp., n[:,1]=unstable
\end{lstlisting}
\begin{itemize}
\item[\textbf{L315}] $T$, the simulation horizon.
\item[\textbf{L316}] Trailing zero appended to \texttt{g\_full}.
\item[\textbf{L317--L318}] Slices the length-$(T+1)$ \texttt{kg\_full} into
``now'' and ``next'' views.
\item[\textbf{L319}] $N_k[t]=\Phi_{kg}\hat g_t+\Phi_{k,kg}\hat k^G_t$.
\item[\textbf{L320--L321}] $N_c[t]$, using both contemporaneous and
one-period-ahead forcing, per \eqref{eq:M-system}.
\item[\textbf{L322}] Stacks $N_k,N_c$.
\item[\textbf{L323}] Rotates into the eigenbasis.
\end{itemize}

\begin{lstlisting}[firstnumber=325]
zeta2 = np.zeros(T)
running = 0.0
for t in range(T - 1, -1, -1):
    running = (n[t, 1] + running) / lam_u
    zeta2[t] = -running
\end{lstlisting}
\begin{itemize}
\item[\textbf{L325--L329}] Solves the unstable (jump) component
\emph{backward} (a naive forward iteration would blow up since
$|\lambda_u|>1$), giving the discounted present value of all remaining
forcing: $\zeta_{2,t}=-\sum_{j=0}^{T-1-t}n_{t+j,1}/\lambda_u^{j+1}$.
\end{itemize}

\begin{lstlisting}[firstnumber=331]
zeta1 = np.zeros(T)
zeta1[0] = (k0 - V[0, 1] * zeta2[0]) / V[0, 0]
for t in range(T - 1):
    zeta1[t + 1] = lam_s * zeta1[t] + n[t, 0]

z = V @ np.vstack([zeta1, zeta2])
return z[0, :], z[1, :]
\end{lstlisting}
\begin{itemize}
\item[\textbf{L331--L332}] Pins $\zeta_{1,0}=(k_0-V_{01}\zeta_{2,0})/V_{00}$
from the true initial condition $\hat k_0=k_0$.
\item[\textbf{L333--L334}] Propagates $\zeta_{1,t}$ forward (safe since
$|\lambda_s|<1$).
\item[\textbf{L336}] Rotates back: $\begin{pmatrix}\hat k_t\\\hat
c_t\end{pmatrix}=V\zeta_t$.
\item[\textbf{L337}] Returns $\{\hat k_t\},\{\hat c_t\}$.
\end{itemize}

\section{\texttt{simulate\_transition}: assembling the full path (L340--L399)}
\label{sec:transition}

\begin{lstlisting}[firstnumber=340]
@dataclass
class TransitionPath:
    years: np.ndarray
    Y: np.ndarray  # % deviation from the steady state being linearized around
    C: np.ndarray
    I: np.ndarray
    K: np.ndarray
    KG: np.ndarray
    N: np.ndarray
    W: np.ndarray
    G: np.ndarray
    r_bp: np.ndarray  # real interest rate, deviation in basis points (annualized)
\end{lstlisting}
\begin{itemize}
\item[\textbf{L340--L351}] Holds one array per series shown in the app's
``Transition dynamics'' charts, all in percent deviations from a reference
steady state.
\end{itemize}

\begin{lstlisting}[firstnumber=354]
def simulate_transition(ss_ref: SteadyState, g_path: np.ndarray,
                         k0: float = 0.0, kg0: float = 0.0, T_sim: int = 300) -> TransitionPath:
    coeffs = _log_linear_coeffs(ss_ref)
    y_k, y_kg, y_c, y_g, s_GR = coeffs["y_k"], coeffs["y_kg"], coeffs["y_c"], coeffs["y_g"], coeffs["s_GR"]

    g_full = np.zeros(T_sim)
    n_g = min(len(g_path), T_sim)
    g_full[:n_g] = g_path[:n_g]

    L = ss_ref.L
    tau = ss_ref.tau
\end{lstlisting}
\begin{itemize}
\item[\textbf{L354--L355}] Signature: steady state to linearize around,
government-spending shock path, initial capital log-deviations
($k_0,kg_0$, default $0$), simulation horizon.
\item[\textbf{L359--L360}] Computes and unpacks the coefficients.
\item[\textbf{L362--L364}] Zero-pads \texttt{g\_path} to length
$T_{\text{sim}}$.
\item[\textbf{L366--L367}] Unpacks $L,\tau$.
\end{itemize}

\begin{lstlisting}[firstnumber=369]
kg_full_ext = _kg_hat_path(ss_ref.delta, g_full, kg0)  # length T_sim+1
k, c = _solve_saddle_path(coeffs["Phi_kk"], coeffs["Phi_kc"], coeffs["Phi_kg"], coeffs["Phi_k_kg"],
                           coeffs["coef_c"], coeffs["coef_k"], coeffs["coef_g"], coeffs["coef_kg"],
                           k0, g_full, kg_full_ext)
kg_full = kg_full_ext[:T_sim]
y = y_k * k + y_kg * kg_full + y_c * c + y_g * g_full
\end{lstlisting}
\begin{itemize}
\item[\textbf{L369}] Solves public capital's path first (it does not
depend on $k,c$).
\item[\textbf{L370--L372}] Calls the saddle-path solver.
\item[\textbf{L373}] Trims the extra period.
\item[\textbf{L374}] $\hat y_t=y_k\hat k_t+y_{kg}\hat k^G_t+y_c\hat
c_t+y_g\hat g_t$, the static block \eqref{eq:static}.
\end{itemize}

\begin{lstlisting}[firstnumber=376]
if ss_ref.distortionary:
    n = L * (y - c) - (L * tau / (1.0 - tau)) * (g_full - y)
else:
    n = L * (y - c)
w = y - n  # wage = MPL, log-deviation

i = (y - ss_ref.s_C * c - s_GR * g_full) / ss_ref.s_I
\end{lstlisting}
\begin{itemize}
\item[\textbf{L376--L379}] Implements \eqref{eq:n-distortionary} or its
lump-sum simplification directly.
\item[\textbf{L380}] $\hat w_t=\hat y_t-\hat n_t$.
\item[\textbf{L382}] $\hat\imath_t=(\hat y_t-s_C\hat c_t-s_{IG}\hat
g_t)/s_I$, from Section~\ref{sec:capaccumsec}.
\end{itemize}

\begin{lstlisting}[firstnumber=384]
c_pct = 100.0 * c
y_pct = 100.0 * y
i_pct = 100.0 * i
k_pct = 100.0 * k
kg_pct = 100.0 * kg_full
n_pct = 100.0 * n
w_pct = 100.0 * w
g_pct = 100.0 * g_full

r_dev_bp = np.zeros_like(c)
r_dev_bp[:-1] = 10000.0 * (1.0 + ss_ref.r) * (c[1:] - c[:-1])
r_dev_bp[-1] = r_dev_bp[-2]

years = np.arange(T_sim)
return TransitionPath(years=years, Y=y_pct, C=c_pct, I=i_pct, K=k_pct, KG=kg_pct, N=n_pct,
                       W=w_pct, G=g_pct, r_bp=r_dev_bp)
\end{lstlisting}
\begin{itemize}
\item[\textbf{L384--L391}] Converts every log-deviation to percent.
\item[\textbf{L393--L395}] Real interest rate deviation in basis points,
via a forward difference of $\{\hat c_t\}$ ($dr_{t+1}\approx(1+r)(\hat
c_{t+1}-\hat c_t)$); the last period repeats the second-to-last value.
\item[\textbf{L397--L399}] Assembles and returns the \texttt{TransitionPath}.
\end{itemize}

\section{\texttt{run\_experiment}: a full policy comparison (L406--L476)}
\label{sec:expsec}

\begin{lstlisting}[firstnumber=406]
@dataclass
class PolicyExperiment:
    ss_old: SteadyState
    ss_new: SteadyState
    path: Optional[TransitionPath]
    multiplier_long_run: float
    multiplier_impact: Optional[float]
\end{lstlisting}
\begin{itemize}
\item[\textbf{L406--L412}] Bundles the before/after steady states, the
transition path, and the two headline multipliers.
\end{itemize}

\begin{lstlisting}[firstnumber=415]
def run_experiment(theta_N: float, delta: float, r: float, A: float, theta_G: float,
                    s_G_old: float, s_G_new: float, f_IG: float, f_GT: float,
                    N_target: float, financing: Financing,
                    permanent: bool, duration_years: int = 4, T_sim: int = 300) -> PolicyExperiment:
    distortionary = financing == "income_tax"
    s_IG_old, s_GT_old = s_G_old * f_IG, s_G_old * f_GT
    theta_L = calibrate_theta_L(theta_N, delta, r, A, theta_G, s_IG_old, s_GT_old,
                                 distortionary, N_target)

    ss_old = steady_state_for_policy(theta_N, delta, r, A, theta_G, s_G_old, f_IG, f_GT,
                                       theta_L, distortionary)
    ss_new = steady_state_for_policy(theta_N, delta, r, A, theta_G, s_G_new, f_IG, f_GT,
                                       theta_L, distortionary)
\end{lstlisting}
\begin{itemize}
\item[\textbf{L415--L418}] Signature: technology/preferences, $\theta_G$,
old/new \emph{total} government-spending ratios (the ``$\Delta$G/Y'' experiment),
composition fractions (held fixed across the experiment), labor-supply
target, financing rule, shock duration.
\item[\textbf{L422}] Translates \texttt{financing} into the boolean
\texttt{distortionary}.
\item[\textbf{L423}] Splits the \emph{baseline} total into its two
composition levels, for calibrating $\theta_L$.
\item[\textbf{L424--L425}] Calibrates $\theta_L$ once, at the baseline
policy --- used for \emph{both} steady states below.
\item[\textbf{L427--L430}] Builds both steady states, holding
$(f_{IG},f_{GT})$ and $\theta_L$ fixed, varying only the total.
\end{itemize}

\begin{notebox}
\textbf{On what ``1. Comparative steady states'' compares.} The app's own
Panel 1 does \emph{not} directly display \texttt{ss\_old}/\texttt{ss\_new}
from this function. Since \texttt{run\_experiment} recalibrates $\theta_L$
fresh from whatever the \emph{current} sidebar baseline happens to be, its
$N$ always lands on exactly \texttt{N\_target} by construction --- which
would silently mask any real labor-supply response to composition,
financing, $\theta_G$, or $r$ changes (dumping the whole adjustment onto
consumption instead). The app therefore computes a separate ``current
settings'' steady state for Panel 1 using \texttt{steady\_state\_for\_policy}
directly, with a theta\_L held \emph{fixed} at a benchmark calibration
shared with the comparison's left-hand column, so $N$ moves honestly there.
\texttt{run\_experiment}'s own recalibration remains exactly right for its
actual purpose here: isolating the effect of the $\Delta$G/Y lever specifically,
holding every other structural parameter (including preferences) fixed on
both sides of that one comparison.
\end{notebox}

\begin{lstlisting}[firstnumber=432]
dY = ss_new.Y - ss_old.Y
dG = ss_new.G - ss_old.G
multiplier_long_run = dY / dG if abs(dG) > 1e-12 else float("nan")
\end{lstlisting}
\begin{itemize}
\item[\textbf{L432--L434}] The long-run output multiplier, $\Delta Y/\Delta
G$, an \emph{exact} finite difference of two closed-form steady states ---
not affected by the log-linear machinery at all.
\end{itemize}

\begin{lstlisting}[firstnumber=436]
if permanent:
    ss_ref = ss_new
    k0 = np.log(ss_old.K / ss_new.K)
    kg0 = np.log(ss_old.KG / ss_new.KG) if ss_new.KG > 0 else 0.0
    g_path = np.zeros(1)
else:
    ss_ref = ss_old
    k0 = 0.0
    kg0 = 0.0
    dG_frac = (ss_new.G - ss_old.G) / ss_old.G if ss_old.G != 0 else 0.0
    g_level = np.log(1.0 + dG_frac)
    g_path = np.concatenate([np.full(duration_years, g_level), np.zeros(1)])
\end{lstlisting}
\begin{itemize}
\item[\textbf{L436}] \textbf{Permanent:} linearize around $\text{ss\_ref}=\text{ss\_new}$.
\item[\textbf{L437}] $\hat k_0=\log(K_{\text{old}}/K_{\text{new}})$ ---
capital is predetermined, still at its old level.
\item[\textbf{L438}] Same logic for public capital, guarded against
$K^G_{\text{new}}=0$.
\item[\textbf{L439}] No additional shock beyond what's baked into
\texttt{ss\_ref}.
\item[\textbf{L441}] \textbf{Temporary:} linearize around the common
$\text{ss\_ref}=\text{ss\_old}$.
\item[\textbf{L443--L444}] Both initial gaps are $0$.
\item[\textbf{L445--L447}] A finite-duration government-spending shock
path.
\end{itemize}

\begin{lstlisting}[firstnumber=449]
path = simulate_transition(ss_ref, g_path, k0=k0, kg0=kg0, T_sim=T_sim)

Y_level_0 = ss_ref.Y * float(np.exp(path.Y[0] / 100.0))
dY0 = Y_level_0 - ss_old.Y
dG0 = ss_new.G - ss_old.G
multiplier_impact = dY0 / dG0 if abs(dG0) > 1e-12 else float("nan")
\end{lstlisting}
\begin{itemize}
\item[\textbf{L449}] Calls Section~\ref{sec:transition}'s function.
\item[\textbf{L451}] Converts the impact-period deviation back into a
level.
\item[\textbf{L452--L454}] The impact multiplier, always relative to the
original steady state.
\end{itemize}

\begin{lstlisting}[firstnumber=456]
def shift(field_ref, X_ref, X_old):
    return field_ref + 100.0 * np.log(X_ref / X_old)

path.Y = shift(path.Y, ss_ref.Y, ss_old.Y)
path.C = shift(path.C, ss_ref.C, ss_old.C)
path.I = shift(path.I, ss_ref.I, ss_old.I)
path.K = shift(path.K, ss_ref.K, ss_old.K)
if ss_ref.KG > 0 and ss_old.KG > 0:
    path.KG = shift(path.KG, ss_ref.KG, ss_old.KG)
else:
    path.KG = np.zeros_like(path.KG)
path.N = shift(path.N, ss_ref.N, ss_old.N)
path.W = shift(path.W, ss_ref.w, ss_old.w)
if permanent:
    path.G = np.full_like(path.G, 100.0 * np.log(ss_new.G / ss_old.G))
else:
    path.G = shift(path.G, ss_ref.G, ss_old.G)
\end{lstlisting}
\begin{itemize}
\item[\textbf{L456--L457}] Re-expresses a path from \% deviation from
\texttt{ss\_ref} into \% deviation from \texttt{ss\_old}.
\item[\textbf{L459--L462}] Applied to $Y,C,I,K$, so every chart is measured
relative to the \emph{original} steady state.
\item[\textbf{L463--L466}] Same shift for $K^G$, guarded against
$\log(0/0)$.
\item[\textbf{L467--L468}] Same shift for $N,w$.
\item[\textbf{L469--L472}] $G$: constant line under a permanent shock;
ordinary shift under a temporary one.
\end{itemize}

\begin{lstlisting}[firstnumber=474]
return PolicyExperiment(ss_old=ss_old, ss_new=ss_new, path=path,
                         multiplier_long_run=multiplier_long_run,
                         multiplier_impact=multiplier_impact)
\end{lstlisting}
\begin{itemize}
\item[\textbf{L474--L476}] Assembles and returns the full
\texttt{PolicyExperiment}.
\end{itemize}

\section{\texttt{public\_investment\_long\_run}: replicating the paper's actual Table 4 (L483--L567)}
\label{sec:table4}

This section was rewritten after reading Baxter \& King's Table 4 (their
page 330) directly, and verified line-for-line against the published
numbers. The paper's own note under the table states its exact
methodology, which the code now follows precisely:

\begin{itemize}
\item Public investment is fixed at $s_{I^G}=0.05$ throughout the whole
table --- \textbf{not} tied to the app's live composition slider elsewhere
(which can go to $0\%$).
\item ``In each case, shifts in purchases are financed via lump-sum
taxation'' --- Table 4 is \emph{always} lump-sum, regardless of the
sidebar's financing toggle.
\item At $\theta_G=0$, the paper's own row reads $\Delta Y/\Delta I^G=1.16$,
\emph{not} $0$: ``the first row of the table replicates the results for
basic government purchases obtained in Section III above'' (the ordinary
lump-sum multiplier). There is no productivity distinction between $I^G$
and other government spending at $\theta_G=0$, so a marginal dollar of
either has the identical negative-wealth-effect labor-supply response.
\end{itemize}

\begin{notebox}
\textbf{Why a small, live-slider-driven $s_{I^G}$ blew up.} An earlier
version of this function set $s_{I^G}=\max(f_{IG}\cdot s_G,\,10^{-4})$,
tied to the sidebar's composition slider, and perturbed it by
$h=10^{-4}$ for the finite difference. When the slider was at $0\%$, this
forced $s_{I^G}-h$ to land at (or below) $0$ exactly. But $K^G$ is an
\emph{essential} input in the production function whenever $\theta_G>0$
(it enters multiplicatively, $Y\propto(K^G)^{\theta_G}$): as $K^G\to0$,
$Y\to0$ \emph{regardless of $K,N$}. The two finite-difference evaluation
points ($s_{I^G}\mp h$) straddled this singularity, producing wildly
different $\alpha$ values on each side and a finite difference that
exploded to values in the billions for large $\theta_G$. The fix
(L489-524) hardcodes $s_{I^G}=0.05$ (safely away from $0$) independent of
any sidebar control, exactly matching the paper's own fixed calibration
choice.
\end{notebox}

\begin{lstlisting}[firstnumber=483]
def public_investment_long_run(theta_N: float, delta: float, r: float, A: float,
                                 s_other: float, N_target: float,
                                 theta_G_grid: np.ndarray, s_IG: float = 0.05) -> "dict[str, np.ndarray]":
\end{lstlisting}
\begin{itemize}
\item[\textbf{L483--L485}] Signature: \texttt{s\_other} is other,
always-unproductive, resource-using government spending held fixed ---
what used to be the separate ``basic purchases'' $G_B$ category before it
was folded into $I^G$ in the main model (Section~\ref{sec:supply}).
Reintroducing it \emph{here specifically} (not in the main model) is what
lets this table replicate the paper's actual setup: a baseline total
government share (matching the sidebar's $G/Y$) of which only $s_{I^G}=0.05$
is the productive investment under study, and the rest ($s_{\text{other}}$)
remains inert at every $\theta_G$ on the grid. \texttt{N\_target} replaces
the old \texttt{theta\_L} parameter --- $\theta_L$ is now calibrated
\emph{inside} this function (L518--L525) rather than by the caller, since
its calibration is specific to this table's own $(s_{IG},s_{\text{other}})$
split and would otherwise require a second, easy-to-mismatch calibration
call in \texttt{app.py}.
\end{itemize}

\textbf{Deriving \texttt{direct} and \texttt{k\_adj} (L529--L530).} Holding
private capital and labor fixed, net output $Y-I^G$ is maximized when
$s_{I^G}=\theta_G$ (``zero net resource use''); more generally,
\eq{ \frac{dY}{dI^G}\bigg|_{K,N\text{ fixed}} = \frac{\partial Y}{\partial K^G}\cdot\frac{dK^G}{dI^G}
   = \left(\theta_G\frac{Y}{K^G}\right)\cdot\frac{1}{\delta} = \frac{\theta_G}{s_{I^G}} }
(using $K^G\delta=I^G=s_{I^G}Y$ at steady state). \textbf{This is the fix
verified against the paper}: an earlier version of this code used
\texttt{direct = theta\_G\_grid} directly (i.e.\ $\theta_G$ alone, implicitly
assuming $s_{I^G}=1$), which does not match the paper's reported ratios at
all (e.g.\ at $\theta_G=0.05$ the paper reports direct effect $=1.00$, not
$0.05$). With $1/(1-\theta_K)$ the standard capital-deepening amplification
factor for letting $K$ (but not $N$) adjust:
\eq{ \texttt{direct} = \frac{\theta_G}{s_{I^G}}, \qquad \texttt{k\_adj} = \frac{\texttt{direct}}{1-\theta_K}. }
Both are now verified to match the paper's Table 4 exactly (e.g.\
$\theta_G=0.05\Rightarrow\text{direct}=1.00,\ \text{k\_adj}=1.72$;
$\theta_G=0.40\Rightarrow\text{direct}=8.00,\ \text{k\_adj}=13.79$).

\begin{lstlisting}[firstnumber=516]
theta_K = 1.0 - theta_N

tau0 = s_IG + s_other
supply0 = _supply_side_impl(theta_N, delta, r, A, tau0, 0.0, s_IG, False)
s_C0 = 1.0 - supply0["s_I"] - s_IG - s_other
if s_C0 <= 0:
    raise ValueError("Government share too large: steady-state consumption share <= 0.")
theta_L = theta_N * (1.0 - N_target) / (s_C0 * N_target)  # tax_factor=1 always (lump-sum)
if theta_L <= 0:
    raise ValueError("Implied theta_L <= 0; adjust N_target or the spending settings.")

direct = theta_G_grid / s_IG
k_adj = direct / (1.0 - theta_K)
\end{lstlisting}
\begin{itemize}
\item[\textbf{L516}] $\theta_K=1-\theta_N$, computed directly rather than
via a full \texttt{steady\_state} call (this function never needs a
\texttt{SteadyState} object, only $\theta_K$).
\item[\textbf{L518--L521}] Calibrates $\theta_L$ once, at $\theta_G=0$
(the ``direct effect'' baseline, always lump-sum so \texttt{tax\_factor}$=1$),
so labor hits \texttt{N\_target} under the total baseline share
$\tau_0=s_{I^G}+s_{\text{other}}$ --- an inlined version of
\texttt{calibrate\_theta\_L} (Section~\ref{sec:calib}) specialized to
this table's $s_C=1-s_I-s_{IG}-s_{\text{other}}$ formula (which the main
\texttt{calibrate\_theta\_L} cannot express, since it has no
$s_{\text{other}}$ concept).
\item[\textbf{L523--L524}] Guards against infeasible calibration.
\item[\textbf{L529--L530}] Implements the closed-form \texttt{direct} and
\texttt{k\_adj} formulas above.
\end{itemize}

\begin{lstlisting}[firstnumber=532]
both = np.zeros_like(theta_G_grid)
dC = np.zeros_like(theta_G_grid)
dI = np.zeros_like(theta_G_grid)
for idx, theta_G in enumerate(theta_G_grid):
    h = min(1e-4, s_IG * 0.05)  # keep the perturbation well inside (0, s_IG)
    m1 = _public_capital_output(theta_N, theta_K, delta, r, A, s_other, theta_L,
                                 theta_G, s_IG - h)
    m2 = _public_capital_output(theta_N, theta_K, delta, r, A, s_other, theta_L,
                                 theta_G, s_IG + h)
    both[idx] = (m2["Y"] - m1["Y"]) / (m2["IG"] - m1["IG"])
    dC[idx] = (m2["C"] - m1["C"]) / (m2["IG"] - m1["IG"])
    dI[idx] = (m2["I"] - m1["I"]) / (m2["IG"] - m1["IG"])

return dict(theta_G=theta_G_grid, direct=direct, k_adj=k_adj, both=both, dC=dC, dI=dI)
\end{lstlisting}
\begin{itemize}
\item[\textbf{L536}] $h=\min(10^{-4},\,0.05\,s_{I^G})$ --- with
$s_{I^G}=0.05$ fixed, this evaluates to $h=2.5\times10^{-4}$ by default,
but the $0.05\,s_{I^G}$ cap is a defensive measure: it guarantees the
perturbation always stays a small \emph{fraction} of $s_{I^G}$ itself
(rather than a fixed absolute step that could exceed $s_{I^G}$ if this
function were ever called with a smaller custom $s_{I^G}$), keeping both
evaluation points $s_{I^G}\mp h$ safely on the same side of the
$K^G\to0$ singularity discussed above.
\item[\textbf{L537--L543}] \textbf{Column (iv--vi), full general
equilibrium} (both $K$ and $N$ adjust): a \textbf{numerical derivative} ---
the exact nonlinear steady state at $s_{I^G}\mp h$, forming the centered
finite difference $\Delta Y/\Delta I^G\approx(Y_2-Y_1)/(I^G_2-I^G_1)$ (and
likewise for $C,I$). Built from two calls to the closed-form
\texttt{\_public\_capital\_output} solver, so it never suffers a
root-finder's convergence failure --- only the (now-fixed) evaluation-point
singularity discussed above.
\item[\textbf{L545}] Returns all five series plus the grid, for the app to
tabulate.
\end{itemize}

\begin{lstlisting}[firstnumber=548]
def _public_capital_output(theta_N, theta_K, delta, r, A, s_other, theta_L, theta_G, s_IG):
    tau = s_IG + s_other
    supply = _supply_side_impl(theta_N, delta, r, A, tau, theta_G, s_IG, False)
    s_I = supply["s_I"]
    s_C = 1.0 - s_I - s_IG - s_other
    if s_C <= 0:
        s_C = 1e-6
    N = theta_N / (theta_L * s_C + theta_N)
    Y = supply["alpha"] * N
    K = supply["kappa"] * N
    KG = supply["KG"] * N
    C = s_C * Y
    I = s_I * Y
    IG = s_IG * Y
    return dict(Y=Y, C=C, I=I, IG=IG, K=K, KG=KG, N=N)
\end{lstlisting}
\begin{itemize}
\item[\textbf{L548}] Signature: a lean, single-purpose steady-state solve
(deliberately not routed through the general \texttt{steady\_state}
function, both to avoid dataclass-construction overhead in this
per-grid-point, per-perturbation inner loop, and because it needs the
$s_{\text{other}}$ resource-constraint term that \texttt{steady\_state}
cannot express). \textbf{Always lump-sum} (\texttt{distortionary} is no
longer a parameter at all --- hardcoded to \texttt{False} throughout, since
the paper's Table 4 always is, independent of the sidebar's financing
toggle elsewhere in the app).
\item[\textbf{L549}] $\tau=s_{IG}+s_{\text{other}}$, the total government
share for this cross-section.
\item[\textbf{L550}] Calls Section~\ref{sec:supply}'s price system.
\item[\textbf{L552--L554}] $s_C=1-s_I-s_{IG}-s_{\text{other}}$ --- both
resource-using categories reduce the consumption share; floored at a tiny
positive number rather than raising, so the finite-difference step in
\texttt{public\_investment\_long\_run} never crashes even at the edge of
feasibility.
\item[\textbf{L556}] $N=\theta_N/(\theta_L s_C+\theta_N)$ --- the lump-sum
labor-supply formula (no \texttt{tax\_factor}, since it is always $1$
here), identical in form to \texttt{steady\_state}'s L171 with
$\text{tax\_factor}=1$.
\item[\textbf{L557--L562}] Computes and returns $Y,K,K^G,C,I,I^G,N$.
\end{itemize}

\end{document}
